\section{Full Ablation Tables}
\label{sec:app-ablations}

\FloatBarrier
\subsection{Input Representation Comparison (E02)}

\begin{table}[htbp]
\centering
\caption{Representation comparison on \speck 5r (5 seeds).
         R2\_xor\_diff is our default.}
\label{tab:repr-speck}
\begin{tabular}{@{}lcc@{}}
\toprule
\textbf{Representation} & \textbf{Accuracy (\%)} & \textbf{Dim.} \\
\midrule
R2\_xor\_diff  & $88.65 \pm 0.37$ & 96 \\
R8\_statistical & $88.62 \pm 0.36$ & 96 \\
R4\_bit\_sliced & $86.53 \pm 1.26$ & 96 \\
R1\_raw\_pair   & $86.00 \pm 1.51$ & 64 \\
R3\_concat      & $86.01 \pm 1.86$ & 64 \\
R5\_word\_level & $59.80 \pm 1.60$ & 6 \\
\bottomrule
\end{tabular}
\end{table}

\begin{table}[htbp]
\centering
\caption{Representation comparison across cipher families.
         R2\_xor\_diff is optimal for SPECK and PRESENT but not SIMON.}
\label{tab:repr-cross}
\begin{tabular}{@{}lcccc@{}}
\toprule
\textbf{Repr.} & \textbf{SPECK 5r} & \textbf{SIMON 7r} & \textbf{PRESENT 5r} \\
\midrule
R1\_raw        & 86.00 & \textbf{85.58} & 66.13 \\
R2\_xor\_diff  & \textbf{88.65} & 81.35 & \textbf{76.09} \\
R3\_concat     & 86.01 & 85.46 & 66.22 \\
R4\_bit\_sliced & 86.53 & 85.63 & 66.56 \\
R5\_word       & 59.80 & 60.74 & 63.79 \\
R8\_statistical & 88.62 & 81.48 & 76.11 \\
\bottomrule
\end{tabular}
\end{table}

Key finding: R2\_xor\_diff is not universally optimal.
For \simon, raw and bit-sliced representations outperform xor-diff by 4.3\pp, likely because \simon's AND-rotation nonlinearity creates features better captured by raw bit patterns than by XOR differences.
We retain R2\_xor\_diff across families for controlled comparison.

\FloatBarrier
\subsection{Supplementary Representations (E02 Supplement)}

\begin{table}[htbp]
\centering
\caption{Supplementary representation variants on \speck 5r (5 seeds).}
\label{tab:repr-supplement}
\begin{tabular}{@{}lcc@{}}
\toprule
\textbf{Variant} & \textbf{Accuracy (\%)} & \textbf{Note} \\
\midrule
R7\_sequential     & $100.00 \pm 0.00$ & White-box (all round outputs) \\
R6\_joint\_pc      & $99.98 \pm 0.00$  & Joint plaintext-ciphertext \\
R9\_noise\_only    & $87.22 \pm 0.52$  & Noise-augmented \\
R9\_both           & $85.74 \pm 0.52$  & Noise + mask \\
R9\_mask\_only     & $84.71 \pm 0.61$  & Mask-augmented \\
\bottomrule
\end{tabular}
\end{table}

R7\_sequential (100\%) provides all intermediate round outputs and serves as a white-box sanity check.
R6\_joint\_pc (99.98\%) includes plaintext bits, confirming that plaintext access trivializes the task.

\FloatBarrier
\subsection{Data Efficiency (E14)}

\begin{table}[htbp]
\centering
\caption{Accuracy vs.\ training set size on \speck 5r.}
\label{tab:data-efficiency}
\begin{tabular}{@{}rcc@{}}
\toprule
\textbf{Samples} & \textbf{Accuracy (\%)} \\
\midrule
1,000     & $74.47 \pm 1.23$ \\
5,000     & $78.78 \pm 2.35$ \\
10,000    & $81.06 \pm 1.05$ \\
50,000    & $84.42 \pm 0.63$ \\
100,000   & $84.73 \pm 0.45$ \\
500,000   & $87.16 \pm 0.53$ \\
1,000,000 & $88.21 \pm 0.53$ \\
\bottomrule
\end{tabular}
\end{table}

Returns diminish beyond 100K samples: doubling from 500K to 1M yields only 1.05\pp.
At 1K samples the MLP already achieves 74.5\%, suggesting neural distinguishers are viable even in data-constrained settings.

\FloatBarrier
\subsection{ResNet vs.\ MLP: Full Comparison}

See \Cref{tab:resnet} in the main text.
The maximum gap is 2.64\pp (\present 5r), and the vulnerability ordering is preserved under both architectures.
On \present at 3--4 rounds, MLP matches or slightly exceeds ResNet (gap: $-0.02$ and $-0.19$\pp respectively), suggesting the deeper architecture provides no benefit when the signal is already saturated.
